\section{Reproducibility and Decoder Regimes}

Reference runs use a separate artifact namespace and freeze execution revision,
model revision, tokenizer, recipe, dataset hash, generation settings, and seed
registry before scientific execution. Snapshots remain attached to the original
reference trajectory bank. Each shard writes atomic item completion records;
a single writer verifies expected shards, row counts, identifiers, fingerprints,
and hashes before aggregation. Only sealed compact bundles may populate the
main tables. Raw trajectories remain outside version control.

\begin{table}[h]
\caption{Distinct decoder regimes and reporting units. The reference row is the
planned primary evidence; the measurement row is preserved legacy evidence.}
\label{tab:decoder-regimes}
\centering\small
\begin{tabular}{lll}
\toprule
Regime & Baseline & Failed denominator \\
\midrule
Reference & One trajectory/item; pass@1 & Sampled failed trajectories \\
Measurement V2 & Eight trajectories/item; pass@8 & Separately identified \\
\bottomrule
\end{tabular}
\end{table}

Measurement-decoder BBH and MBPP are not scheduled for reference full reruns in
this design. Legacy successful runs require their own provenance seal before
secondary integration; an unsealed completion is not evidence approval.
The eventual decoder comparison must record both recipes and use matching
denominators for any directly compared metric.
